\subsection{An unrestricted lift with no bottom supported zeros}\label{sec:unrestricted}
We now prove that the ideal hypothesis marks a genuine boundary for the supported-zero conclusion. This construction changes the zero pattern of coefficients whose traces vanish.

Write $\Gamma=\{\gamma_w^r:r=1,2,\ w\in\{1,2,12,21\}\}$ and $B=\{1,\Sigma\}\cup\Gamma$. Give these basis elements heights
\[
 h_1=0,\qquad h_\Sigma=2,\qquad h_{\gamma_w^r}=|w|-1.
\]
For $x,y\ne1$, let $b(x,y)$ be the largest exponent in the trace coefficient $c_{xy}^{\Sigma}$. All these coefficients are positive throughout the interval $[-b(x,y),b(x,y)]$. The multiplication rule gives
\[
\begin{array}{c|c}
(x,y)&b(x,y)\\ \hline
(\Sigma,\Sigma)&4\\
(\Sigma,\gamma_w^1)\text{ or }(\gamma_w^1,\Sigma)&2+h_{\gamma_w^1}\\
(\Sigma,\gamma_w^2)\text{ or }(\gamma_w^2,\Sigma)&3+h_{\gamma_w^2}\\
(\gamma_u^1,\gamma_v^1)&h_{\gamma_u^1}+h_{\gamma_v^1}+1-\beta(u,v)\\
(\gamma_u^2,\gamma_v^2)&h_{\gamma_u^2}+h_{\gamma_v^2}+1\\
(\gamma_u^1,\gamma_v^2)\text{ or the reverse}&h_x+h_y+1
\end{array}
\]
Here $\beta(u,v)\in\{0,1\}$ is the alternating-triple exponent in Lemma~\ref{lem:table}; it need not be symmetric in $u,v$. Positivity on each entire interval follows from the displayed $W_0,U_i,V_i,G_i,S$: the nontrivial short factors $5D-4$ and $2D^2-5D+4$ have sequences $(5,6,5)$ and $(2,3,6,3,2)$, and the other factors are products of $D$ and $D-1$ with positive constants.

Consequently, for all nonunit $x,y$,
\begin{equation}\label{eq:bottom-radii}
 b(x,y)\le1+h_x+h_y,\qquad
 \max_{z\in\Gamma}\bigl(b(z,y)-h_z\bigr)=1+h_y
 =\max_{z\in\Gamma}\bigl(b(y,z)-h_z\bigr).
\end{equation}
The maxima are attained by choosing $z$ in the opposite cell from $y$; if $y=\Sigma$, choose $z$ in the second cell.

Define lifted coefficients $\widetilde c_{xy}^v$ as follows. Unit products have their exact Kronecker coefficients. For $x,y\ne1$, the coefficient at $\Sigma$ has exactly the support of $c_{xy}^{\Sigma}$ and the same integer amplitudes. At every other output $v\ne\Sigma$, use the support interval
\[
 [-r(x,y;v),r(x,y;v)],\qquad
 r(x,y;v)=3+h_x+h_y-h_v,
\]
and place the trace coefficient $c_{xy}^v$ on it, using old zeros at all remaining weights. These radii are nonnegative. The interval contains the trace support: the only nonzero nonbottom coefficients of nonunit products have the form
$\delta_iD^\alpha\lambda_i^\beta$, of degree at most $2+\alpha+2\beta$, while the word-length identity gives
\[
 3+h_x+h_y-h_v=2+\alpha+2\beta.
\]
Thus this is a well-defined coefficientwise lift to $\mathcal S$.

\begin{proof}[Proof of Theorem~\ref{thm:no-bottom-main}]
Every associativity equation holds on integer amplitudes because the trace algebra is associative. It remains to check the finite occupied supports. A sum of centered intervals is their union, hence has the largest radius, and their Minkowski product adds radii.

A triple involving $1$ is associative by the exact unit coefficients. Let $x,y,z\ne1$. For a nonbottom output $v\ne\Sigma$, every intermediate element $u\in\Gamma$ contributes the same radius to $(xy)z$:
\[
 r(x,y;u)+r(u,z;v)=6+h_x+h_y+h_z-h_v.
\]
The intermediate bottom element contributes at most that radius, since
\[
 b(x,y)+r(\Sigma,z;v)
 \le(1+h_x+h_y)+(5+h_z-h_v).
\]
An intermediate unit can contribute only when $v=z$, and then its radius $3+h_x+h_y$ is smaller. The corresponding calculation for $x(yz)$ has the same largest radius.

At output $\Sigma$, the intermediate cell elements have largest radius
\[
 \max_{u\in\Gamma}\bigl(r(x,y;u)+b(u,z)\bigr)
   =4+h_x+h_y+h_z
\]
by \eqref{eq:bottom-radii}. The intermediate bottom contribution satisfies
\[
 b(x,y)+b(\Sigma,z)\le4+h_x+h_y+h_z.
\]
If the intermediate unit contributes, then $z=\Sigma$, and its radius $3+h_x+h_y$ is again smaller. Reversing the association gives the same largest radius by the other equality in \eqref{eq:bottom-radii}. This proves every support associativity equation and therefore associativity over $\mathcal S$.

The bottom supports were defined to contain precisely the weights with positive trace coefficient. They contain no occupied zero coefficient. In contrast, the unit coefficient of $\Sigma^2$ has support $[-7,7]$ and identically zero trace. It is not the absent coefficient $\tau$, so the span of $\Sigma$ is not an ideal. The construction proves every asserted property.
\end{proof}

The lift has $829$ nonabsent coefficient triples and $6{,}167$ supported zero-trace weights, all outside the bottom output. These finite counts follow by expanding the stated intervals and the $132$ trace coefficients; an explicit finite-set check verifies all $10{,}000$ coefficientwise support associativity equations. They are not minimum claims. No effective parity lift of this particular support table has been established here. The example refutes the unrestricted supported-zero assertion for inverse-semiring lifts; it does not refute a statement restricted to effective dimension rings with a preserved bottom ideal.
